%! Characteristics of Quadratic Functions — Unit 3, Lesson 3.0 — Exit Ticket (Answer Key)
\documentclass[10pt]{article}
\usepackage{algebra2-article}
\usepackage{algebra2-key}   % pulls in -boxes; do NOT also load -boxes
\newcommand{\TallMath}[1]{$\displaystyle #1\rule[-1.4em]{0pt}{3.2em}$}
\newcommand{\lgrid}[4]{%
  \draw[step=1, linegray!40, very thin] (#1,#3) grid (#2,#4);
  \draw[->, semithick, charcoal] (#1,0)--(#2,0) node[below right, font=\scriptsize]{$x$};
  \draw[->, semithick, charcoal] (0,#3)--(0,#4) node[left, font=\scriptsize]{$y$};}

\begin{document}

\pageheader{Unit 3, Lesson 3.0}{Exit Ticket --- Answer Key}
\namedateperiod

\vspace{0.10in}

No notes. The graph below is the quadratic $f(x) = x^2 + 2x - 3$. Use it for all questions.
\begin{center}
\begin{tikzpicture}[scale=0.52]
  \lgrid{-4}{2}{-5}{3}
  \draw[navy, dashed, semithick] (-1,-5)--(-1,3);
  \begin{scope}
    \clip (-4,-5) rectangle (2,3);
    \draw[very thick, forest, domain=-3.6:1.6, samples=75, smooth]
      plot (\x,{(\x)*(\x) + 2*(\x) - 3});
  \end{scope}
  \fill[charcoal] (-1,-4) circle (3pt) node[below right, font=\scriptsize]{$(-1,-4)$};
  \fill[charcoal] (0,-3) circle (3pt);
  \fill[charcoal] (-3,0) circle (3pt);
  \fill[charcoal] (1,0) circle (3pt);
\end{tikzpicture}
\end{center}

\begin{enumerate}[leftmargin=1.8em, itemsep=10pt]
  \item \textbf{Read the features.} Give the vertex, the axis of symmetry, the two zeros, and the
        $y$-intercept.
        \ansline{Vertex $(-1,-4)$; \; axis of symmetry $x=-1$; \; zeros $x=-3$ and $x=1$; \; $y$-intercept $(0,-3)$.}
  \item \textbf{Classify \& justify.} Does the vertex give a maximum or a minimum, and what is that
        value? State the range. Then explain how the \emph{sign of $a$} tells you the parabola opens
        that way.
        \ansline{Minimum; minimum value $-4$. Range: $y\ge -4$, i.e.\ $[-4,\infty)$. Since $a=1>0$, the parabola opens upward (a bowl), so the vertex is the lowest point and the outputs rise from there in both directions.}
  \item \textbf{SOL-style multiple choice.} Which statement about $f(x)=x^2+2x-3$ is \emph{true}?
        \begin{enumerate}[label=\Alph*., leftmargin=1.9em, itemsep=1pt]
          \item $f$ is increasing for all $x$.
          \item $f(x)<0$ only on the interval $-3<x<1$. \textcolor{keyred}{\textbf{$\leftarrow$ correct}}
          \item The range of $f$ is all real numbers.
          \item As $x\to-\infty$, $f(x)\to-\infty$.
        \end{enumerate}
        \begin{itemize}[leftmargin=1.4em, nosep]
          \item[] \textbf{A} is false --- $f$ decreases on $(-\infty,-1)$, then increases. \textbf{C} is false --- the range is $[-4,\infty)$. \textbf{D} is false --- both arms rise ($a>0$), so $f(x)\to+\infty$. \textbf{B} is correct: the graph is below the $x$-axis exactly between the zeros.
        \end{itemize}
\end{enumerate}

\begin{teachernote}
Sort responses into ``got it / point-vs-value slip (min as a point) / sign-of-$a$ slip / interval
slip.'' The telltale error is item~3D or a range of ``all reals'' --- both mean the student has not
connected $a>0$ to ``opens up, has a minimum, both arms rise.'' If several miss it, reopen Lesson~3.1
by shading above/below the axis and marking the vertex on one parabola.
\end{teachernote}

\end{document}
